\begin{lemma}[Situation 7.2, Subcase 2.2]
Every \noderef{HierarchySituationTwoCaseTwoTwoData} is impossible.
Equivalently, Situation 7.2, Case 2 cannot occur when both
\(\alpha\) and the last-edge color \(\theta=\varphi_n(e_p)\) are
missing at \(y_{p-1}\).
\end{lemma}
\begin{proof}
Let \(\mathcal E\) denote the supplied
\noderef{ExtendedTashkinovTree}, and let the branch fields be the
supplied \noderef{HierarchySituationTwoCaseTwoTwoData}.  Write
\(H=C.\mathit{hierarchy}\), let \(q\) be its last stored index,
and put
\[
 A=\noderef{HierarchyStarredStage}(H,q),\qquad
 \ell=\noderef{HierarchyFullTailLength}(C),
\]
\[
 B=\noderef{HierarchyFullTailPrefix}(C,\ell-2),\qquad
 B^+=\noderef{HierarchyFullTailPrefix}(C,\ell-1),\qquad
 D_{n,q}=\noderef{HierarchyStoredColorsAt}(H,q).              \tag{1}
\]
The nested-stage and terminal-stage fields of
\noderef{TashkinovHierarchy} put \(A\) below \(T\) as an exact
\noderef{TreeSequenceInitialSegment}.  Since the datum gives
\(2\leq\ell\) and \(\ell=|T|-|A|\), expanding the two prefixes in
(1) gives
\[
 A\preccurlyeq B\prec B^+\prec T,
 \quad |B|=|T|-2,\quad |B^+|=|T|-1.                          \tag{2}
\]
All relations in (2) are entrywise: they identify each retained vertex
and each individual edge, including parallel edge identities.  The
three position fields of the datum consequently identify
\(y_{p-1}\) with the last vertex of \(B^+\), \(y_p\) with the last
vertex of \(T\), and \(e_p\) with the unique edge added above \(B^+\).
The final ladder transition used below is always the actual certified
\noderef{TashkinovSeriesStep} of \(\mathcal E\); every replacement ETT
will be supplied with a full
\noderef{CorrespondingExtendedTashkinovTree} certificate, rather than
only an abstract recoloring of \(T\).

Apply \noderef{HierarchyCounterexamplePenultimateElementary} to the
properness, exact one-uncoloredness, complete lower-rung six-pack,
fewer-divider induction, rung-local lexicographic minimality, and
large-set fields of the supplied
\noderef{HierarchyMinimumCounterexampleContext}.  Its conclusion is the
derived fact (6.6)
\[
       \noderef{ColoredVertexSetElementary}
          (G,\varphi_n,V(B^+)).                               \tag{3}
\]
In particular, (3) is not being assumed as an extra datum.

We first reprove equation (9) in this `.2' branch.  If \(\alpha\) were
missing at a vertex of \(B\), that vertex and \(y_{p-1}\) would be
distinct entries of \(B^+\) by (2), both missing \(\alpha\), contrary
to (3).  Hence
\[
                    \alpha\notin\overline\varphi_n(B).       \tag{4}
\]
Since \(A\preccurlyeq B\), (4) also says that \(\alpha\) is not
missing on \(A\).  Let
\[
 D_n=\noderef{TashkinovStoredColors}
   (\mathcal E.\mathit{series}.\mathit{data},
      \operatorname{Fin.last}(\mathcal E.\mathit{rounds})).
\]
Unfolding \noderef{HierarchyStoredColorsAt} at the last index gives
\[
                 D_{n,q}=D_n\setminus\overline\varphi_n(A).
\]
Thus membership of \(\alpha\) in \(D_n\), together with its absence
from the missing set of \(A\), would put it in \(D_{n,q}\), contrary
to the datum.  Therefore
\[
              \alpha\notin\overline\varphi_n(B),\qquad
              \alpha\notin D_n.                              \tag{9}
\]

We also need the edge form of (9).  Apply
\noderef{HierarchyTailMissingColorExclusion} with \(i=\ell-1\),
the elementary prefix (3), and the exact last vertex
\(y_{p-1}\) of \(B^+\).  It excludes \(\alpha\) from every exact
tail edge whose position is less than \(|B^+|\), unless the edge is the
divider-zero SE/RE seed and \(\alpha\in D_n\).  The second assertion
of (9) eliminates that exception.  For the remaining edge \(e_p\), the
attachment field of \noderef{TreeSequence}, together with the exact
last indices, makes \(e_p\) incident with \(y_p\).  It has color
\(\theta\), while \(\alpha\) is missing at \(y_p\); expanding
\noderef{MissingColors} therefore gives \(\alpha\ne\theta\).  Hence
\[
 \varphi_n(T.\mathit{edge}(t))\ne\alpha
       \quad\text{for every }t\text{ with }|A|\leq t<|T|.    \tag{10}
\]
The same elementarity argument as for (4), now using the datum's
\(\theta\)-missingness at \(y_{p-1}\), shows that no vertex of \(B\)
misses \(\theta\).  Thus the totalized
\noderef{HierarchyFirstMissingPosition} of \(\theta\) is exactly the
position of \(y_{p-1}\).                                      \(\tag{11}\)

We split according as \(\theta\in D_{n,q}\).

\textbf{The stored-\(\theta\) branch.}
Assume \(\theta\in D_{n,q}\).  The `twoColors' field of
\noderef{GoodHierarchy} supplies distinct colors
\(\varepsilon_1,\varepsilon_2\) with
\[
 C.\mathit{good}.\Gamma^q_\theta
       =\{\varepsilon_1,\varepsilon_2\}.                    \tag{12}
\]
By \noderef{GoodHierarchyReservedDisjoint}, both reserved colors are
missing on the relevant hierarchy stage (and hence on \(A\)) and are
outside \(D_{n,q}\).  Since \(\theta\) is stored, (12) also gives
\(\varepsilon_i\ne\theta\).

Unfold the avoidance field of \noderef{GoodHierarchy} at the stored
color \(\theta\).  By (11), the cutoff in
\noderef{HierarchySegmentColors} is the vertex \(y_{p-1}\); hence
neither reserved color is used on an exact edge from \(A\) through
\(B^+\).  The only remaining tail edge is \(e_p\), whose color is
\(\theta\ne\varepsilon_i\).  Applying
\noderef{HierarchyTailMissingColorExclusion} to \(\theta\) itself says
that \(\theta\) is absent from the same prefix interval, except
possibly on the exact divider-zero SE/RE seed.  We have therefore proved
the full manuscript assertion
\[
 \begin{split}
 &\varepsilon_1,\varepsilon_2
       \notin\varphi_n\langle T-A\rangle,\\
 &\theta\notin\varphi_n\langle B^+-A\rangle,
   \text{ apart only from the exact seed }f_n=e_1
   \text{ when }q=0, A=T_n, \varphi_n(f_n)=\theta\in D_n.
 \end{split}                                                  \tag{14}
\]

Choose the first \(\varepsilon_1\)-missing vertex
\(v_{\varepsilon_1}\) in \(A\).  Equation (4) says that the first
\(\alpha\)-missing vertex is \(y_{p-1}\).  The two colors are
distinct, since \(\varepsilon_1\) is missing on \(A\) and \(\alpha\)
is not.  Both are missing on \(B^+\).  Apply
\noderef{ChangeLemma} with ordered colors
\((\varepsilon_1,\alpha)\).  The first color is in the required
stage-missing slot.  If its conditional order premise is active, its
first witness lies in \(A\), the first \(\alpha\)-witness is
\(y_{p-1}\), and (14) gives the exact unused interval.  This proves
\[
 K_{\varepsilon_1,\alpha}(v_{\varepsilon_1})
      =K_{\varepsilon_1,\alpha}(y_{p-1}).                    \tag{15}
\]

The component in (15) is disjoint from the component rooted at \(y_p\).
This remains true even if the two first-missing witnesses in (15)
coincide: then both colors are missing at that vertex, and unfolding
\noderef{TwoColorIncidenceGraph} makes its component an isolated vertex,
which cannot contain the distinct \(y_p\).  If they are distinct and
the component through \(y_p\) were the same, the component would contain
three distinct vertices of affected-color incidence degree at most one,
missing respectively \(\varepsilon_1,\alpha,\alpha\).  The finite
spanning path from \noderef{TwoColorComponentSpanningPath} and the
degree bound \noderef{TwoColorIncidenceDegreeAtMostTwo} allow only two
endpoints.  The dichotomy in \noderef{TwoColorComponentsDisjoint}
therefore gives disjointness in either case, and the spanning-path lemma
also supplies the possibly singleton \noderef{TwoColorPathComponent}
certificate for the component at \(y_p\).

Define
\[
 \mu_1=\noderef{TwoColorComponentSwap}
   (G,\varphi_n,\varepsilon_1,\alpha,\operatorname{inl}y_p). \tag{16}
\]
Apply \noderef{StableChangeCasesOneTwo} immediately, before using any
stability consequence, with \(\varepsilon_1\) in the stage-missing
slot, (3), (14)--(16), the component separation, and all six fields of
the minimum context.  Retain its
\noderef{HierarchyWeaklyStable} and \noderef{ModColoring} conclusions.

We check the full-tail transport premises pointwise.  Equation (10) and
(14) say that neither affected color occurs on any exact edge of
\(T-A\), so every such individual edge keeps its value.  At a vertex
of \(B^+\) outside the swapped component no incident edge changes.  A
vertex of \(B^+\) inside it cannot be an endpoint: by (3), every endpoint
missing an affected color is the corresponding witness in the disjoint
component (15).  Properness gives one incident edge of each affected
color, and the swap exchanges them, preserving the complete missing set.
Thus
\[
 \overline\mu_1(u)=\overline\varphi_n(u)\quad(u\in V(B^+)),
 \qquad \mu_1(e)=\varphi_n(e)\quad(e\text{ displayed in }T-A).\tag{17}
\]
Finally, neither affected color lies in \(D_{n,q}\): this is reserved-
stored disjointness for \(\varepsilon_1\) and the datum for \(\alpha\).
Every stored color is consequently unaffected at every vertex, so direct
expansion of \noderef{HierarchyFirstMissingPosition}, including its
terminal fallback, gives equality of all stored first positions.

The three facts in (17) and the stored-position equality are precisely
the separate hypotheses of
\noderef{StableChangeCasesOneTwoFullTail}.  That theorem returns an
actual corresponding ETT on the same exact \(T\), maximum property, a
rebuilt base, a full hierarchy, a \noderef{GoodHierarchy}, and
\noderef{SameGoodHierarchy} under \(\mu_1\).

Package these objects as a genuine
\noderef{HierarchyReplacementPresentation} \(P_1\).  Its counterexample
uses the old \(\varepsilon_1\)-missing witness in \(A\), retained by
component separation, and the new \(\varepsilon_1\)-missing vertex
\(y_p\); the latter follows directly from (16), whether the component at
\(y_p\) was a nontrivial path or an isolated vertex.  These vertices are
distinct by (2), so the new hierarchy is non-elementary.  We fill the six
minimum-context fields separately:
\begin{enumerate}
\item properness is \noderef{TwoColorComponentSwapProper};
\item the definition of the swap fixes `none', so
      \noderef{IsOneEdgeUncolored} at the exact edge \(e_0\) is retained;
\item rung agreement of the corresponding ETT transports the complete
      universally quantified lower-rung six-pack;
\item same-good identifies the divider count, so the fewer-divider
      induction field applies unchanged;
\item the terminal tree is literally \(T\), and same-good identifies
      all stages and divider positions, so unfolding
      \noderef{HierarchyCounterexampleMeasure} gives equality of both
      measure coordinates before the old minimality field is applied;
\item weak stability identifies the missing sets on the terminal ladder
      tree and transports the \(2n+11\) estimate.
\end{enumerate}
Thus \(P_1\) contains an actual replacement context, not merely the
claim that the coloring remains stable.

We now change \(\theta\) and \(\varepsilon_1\) inside \(P_1\).  Apply
\noderef{HierarchyCounterexamplePenultimateElementary} to the six fields
just constructed to derive (3) for \(P_1\).  Equations (12), (14), and
(17), transported by same-good, say that \(\theta\) is stored, first
becomes missing at \(y_{p-1}\), and has reserved pair
\(\{\varepsilon_1,\varepsilon_2\}\), while
\(\varepsilon_1\) is stage-missing and is missing at \(y_p\).
Apply \noderef{ChangeLemma} with ordered colors
\((\varepsilon_1,\theta)\).  Its conditional order and unused-edge
premises are (14), and its Cases 1/2 premise again uses the stage-missing
color \(\varepsilon_1\).  It gives
\[
 K^{\mu_1}_{\varepsilon_1,\theta}(v_{\varepsilon_1})
       =K^{\mu_1}_{\varepsilon_1,\theta}(y_{p-1}).            \tag{18}
\]
As above, split on coincidence of the first-missing witnesses.  In the
coincident case their component is isolated; in the distinct case,
equality with the component through \(y_p\) would give three endpoints
of the spanning path, missing \(\varepsilon_1,\theta,\varepsilon_1\).
The same three cited two-color component results therefore make (18)
disjoint from the component at \(y_p\).

Put
\[
 \mu_2=\noderef{TwoColorComponentSwap}
   (G,\mu_1,\varepsilon_1,\theta,\operatorname{inl}y_p).     \tag{19}
\]
Apply \noderef{StableChangeCasesOneTwo} relative to \(P_1\), before
claiming that any exceptional edge is fixed.  It supplies weak stability
and a mod-coloring for (19).  The endpoint/internal-vertex argument just
used for (17) proves missing-set equality throughout \(B^+\).

We verify the split-tail route with \(X=B^+\).  On the interval
\(B^+-A\), (14) excludes both affected colors, except for the possible
divider-zero SE/RE seed carrying \(\theta\).  In that exception,
\(\theta\in D_n\), the singleton connecting-edge equation identifies
the individual seed, and the controlled incident-edge clause of the
\noderef{StableColoring} contained in the weak stability just obtained
fixes it.  Hence every exact prefix edge from \(A\) to \(B^+\) is fixed.

Inspect the final certified step of the actual
\noderef{TashkinovSeries} in \(P_1\).  In PE it gives a
\(\mu_1\)-\noderef{TashkinovClosure} from the terminal ladder tree; in
SE or RE it gives the individual edge \(f_n\), the literal singleton
connecting-edge equation, an exact
\noderef{TreeSequenceAdjoinEdge} seed \(T_n+f_n\), and a closure from
that seed.  In the positive-divider nonparallel case,
\noderef{HierarchyNonparallelFirstEdge} puts that exact seed below
\(A\); at divider zero the base equality and \(2\leq\ell\) put it below
\(B^+\).  The PE base uses the left part of
\noderef{OrderedTreeUnion}.  Therefore
\noderef{TashkinovAugmentationSequenceRestrict} gives the appropriate
old-color route through \(B^+\).

Transport this route to \(\mu_2\).  Up through \(A\), unfold
\noderef{HierarchyWeaklyStable}: its
\noderef{TreeColoringInvariant}, or its divider-zero combined-stability
clauses and \noderef{CombinedStableClosurePreserved}, preserve every
predecessor witness and displayed edge; in the SE/RE seed case use the
controlled-edge argument of the preceding paragraph.  Above \(A\), the
fixed-edge conclusion and the vertexwise missing-set equality on
\(B^+\) are exactly the two premises of
\noderef{TashkinovAugmentationSequenceColoringTransport}.  This gives
the required \(\mu_2\)-route to \(X\) with the same PE source or exact
SE/RE seed.

Under \(\mu_1\), the suffix \(B^+\to T\) is the one
\noderef{TashkinovAugmentation} adjoining \(e_p\): its color is
\(\theta\), missing at \(y_{p-1}\in V(B^+)\).  Under \(\mu_2\), both
affected colors remain missing on \(X\), at witnesses in the disjoint
component (18).  For the stored-position premise,
\(\varepsilon_1\notin D_{n,q}\), while \(\theta\)'s first position is
\(y_{p-1}\) and missing-set equality on \(B^+\) preserves the whole
candidate set through that witness; all other stored colors are
unaffected.  For the Gamma cutoff, reserved-pair disjointness says that
the only last Gamma pair containing \(\varepsilon_1\) is the pair of
\(\theta\), whose first missing position is exactly \(|B^+|\), while
the stored color \(\theta\) lies in no Gamma pair by
\noderef{GoodHierarchyReservedDisjoint}.

These are every hypothesis of
\noderef{StableChangeCasesOneTwoSplitTail}.  It returns an actual
corresponding same-rung maximum ETT on the exact full tree \(T\), a
rebuilt base and hierarchy, full goodness, and same-good data.  In
particular, (19) changes \(e_p\) from \(\theta\) to
\(\varepsilon_1\), while preserving every missing set on \(B^+\).

Build the resulting replacement presentation \(P_2\) explicitly.  Its
non-elementarity uses \(y_{p-1}\) and \(y_p\), which are now both
missing \(\theta\).  Properness and exact one-uncoloredness use the swap
theorem and the preservation of `none'.  Rung agreement transports the
lower-rung six-pack; same-good transports the divider count and hence
the fewer-divider field; literal equality of \(T\) and equality of all
hierarchy stages give equality of both counterexample-measure coordinates
before minimality is transported; and weak stability on the replacement
terminal ladder tree transports the large-set bound.  Compose the two
corresponding-series certificates entry by entry, use the transitivity
in \noderef{sc2} for each positive stable-coloring slot, and compose the
two stage and Gamma equalities in same-good.  This supplies every field
of an original-relative \noderef{HierarchyReplacementPresentation} and
its complete minimum context.

Finally construct a genuine
\noderef{HierarchySituationTwoCaseOneData} for \(P_2\).  Copy the exact
tail indices, their position equations, non-incidence, full-tail length,
and path-number equation.  Use \(\varepsilon_1\) as the new last-edge
color; it is present at \(y_{p-1}\), since it was missing at its distinct
witness in \(A\), (3) made it present at \(y_{p-1}\), and the two
vertexwise transport equalities preserve that fact.  Use \(\theta\) as
the Case 1 color.  It is missing at both terminal vertices, belongs to
the transported \(D_{n,q}\), and has the transported distinct pair
\(\{\varepsilon_1,\varepsilon_2\}\).  The already Sound
\noderef{HierarchySituationTwoSubcaseOneOne}, applied only to this actual
datum and the actual minimum context of \(P_2\), gives a contradiction.

\textbf{The nonstored-\(\theta\) branch.}
Assume now that \(\theta\notin D_{n,q}\).  We prove the paper's counting
assertion (16) without using the `.1'-specific Case 2 exclusion lemma.
Apply \noderef{NineN} to the present ETT and good hierarchy at
\(i=\ell-2\), using (3) and the large-set field of the context.  Its
available set is
\[
 X=\overline\varphi_n(B)
   \setminus\overline\varphi_n(\mathcal B.\mathit{core}-V(T_n))
   \setminus\varphi_n\langle B-A\rangle,
\]
and \(|X|\geq2n+11\).

If \(|X\setminus(\Gamma^q\cup D_{n,q})|\geq5\), deletion of the colors
of the two exact final edges removes at most two elements.  Choose a
remaining color \(\beta\).  It is missing on \(B\), avoids the base-extra
missing set and every exact edge of \(T-A\), and lies outside both
\(D_{n,q}\) and the full Gamma union.

Suppose instead that this cardinality is at most four.  The second
conclusion of \noderef{NineN} supplies seven injectively indexed stored
colors \(\eta_h\), each missing on \(B\), whose three-element blocks
\[
       Q_h=C.\mathit{good}.\Gamma^q_{\eta_h}\cup\{\eta_h\}
\]
avoid every exact edge of \(B-A\).  These blocks are pairwise disjoint:
the stored colors are distinct, the Gamma pairs are disjoint by
\noderef{GoodHierarchy}, and
\noderef{GoodHierarchyReservedDisjoint} separates every Gamma color from
every stored color.  Consequently each of the two final edge colors meets
at most one block, so one whole block avoids both edges.

We record also why the two Gamma colors of this chosen block avoid the
base-extra missing set.  Follow either Gamma color backwards through the
finite unique-refresh clauses of \noderef{GoodHierarchy}.  At index zero
it is missing on \(T_n\); if it is first refreshed at a positive index,
it is missing on the displayed difference between that stage and the
preceding starred stage.  In either case its supplied missing witness is
outside \(\mathcal B.\mathit{core}-V(T_n)\): the zero witness lies in
\(T_n\), while a positive refresh witness lies beyond a starred stage
which already contains the base core.  Every one of these vertices and
every base-extra vertex lies in the elementary set \(V(B^+)\).  Hence
(3) forbids the same Gamma color from also being missing at a base-extra
vertex.  If the pair was not refreshed at an index, its equality with the
preceding pair continues the finite descent.  This proves the required
base-extra exclusion for both pair members.

Choose \(\beta\) from the two-element Gamma pair in the selected block.
Its stage-missing witness lies in \(A\preccurlyeq B\), reserved-stored
disjointness puts it outside \(D_{n,q}\), and the preceding paragraphs
give all other exclusions.  Combining both cardinality branches proves
exactly
\[
 \begin{split}
 \beta\in{}&\overline\varphi_n(B)
 \setminus\overline\varphi_n(\mathcal B.\mathit{core}-V(T_n))
 \setminus\varphi_n\langle T-A\rangle\setminus D_{n,q},\\
 &\text{and either }\beta\notin\Gamma^q,
 \text{ or }\beta\in\Gamma^q_{\eta_h}
 \text{ for some }\eta_h\in D_{n,q}\cap
       \overline\varphi_n(B).
 \end{split}                                                  \tag{16}
\]

Unfold \noderef{HierarchyStoredColorsAt}.  If \(q>0\) and
\(\beta\in D_n\), its exclusion from \(D_{n,q}\) forces it to be
missing on \(A=T_{n,q}\).  If \(q=0\), the same argument first makes
it missing on the base core; (16) excludes a witness in the base-extra
part, so either it is missing on \(T_n\) or it is outside \(D_n\).
Thus
\[
 \begin{array}{ll}
 q>0:&\beta\in\overline\varphi_n(T_{n,q})\text{ or }\beta\notin D_n,\\
 q=0:&\beta\in\overline\varphi_n(T_n)\text{ or }\beta\notin D_n.
 \end{array}                                                  \tag{17}
\]

Let \(v_\beta\in V(B)\) be the first vertex missing \(\beta\).  By
(4), the first \(\alpha\)-missing vertex is \(y_{p-1}\), so
\(v_\beta\ne y_{p-1}\) and \(\beta\ne\alpha\).  We must distinguish
the exact premise accepted by the Cases 1/2 stable-change theorem from
the weaker disjunction (17).  Write \({\cal C}_{\beta\alpha}\) for
\[
 \begin{split}
 &(q>0\ \hbox{ and }
       \beta\in\overline\varphi_n(T_{n,q}))\quad\hbox{or}\\
 &(q=0\ \hbox{ and }\ [\,
       \beta\in\overline\varphi_n(T_n)\ \hbox{ or }
       (\beta\in\overline\varphi_n(A),
        \beta,\alpha\notin D_n)\,]).                       \tag{18}
 \end{split}
\]
This is word for word the Cases 1/2 proposition in the statement of
\noderef{StableChangeCasesOneTwo}, with the ordered colors
\((\beta,\alpha)\).  Split on \({\cal C}_{\beta\alpha}\).

First apply \noderef{ChangeLemma}, before invoking either stable-change
provider.  Equation (4) says that \(\alpha\) is not missing on \(A\),
so its conditional order-and-unused premise is active.  Exact
prefixhood gives \(v_\beta\prec y_{p-1}\), and (16) excludes \(\beta\)
from every displayed edge from \(A\) to that second witness.  If
\(q>0\), (17) either puts \(\beta\) in the required stage-missing slot,
or (16) and the datum put both \(\beta\) and \(\alpha\) outside
\(D_{n,q}\).  If \(q=0\), (17) either puts \(\beta\) in the
\(T_n\)-missing slot, or (17) and (9) put both colors outside \(D_n\).
These are precisely the first two alternatives in the statement of the
cited change lemma.  It follows that
\[
        K_{\beta,\alpha}(v_\beta)
             =K_{\beta,\alpha}(y_{p-1}).                    \tag{20}
\]

In the complementary stable-change branch, where
\({\cal C}_{\beta\alpha}\) fails, \(\beta\) is not missing on \(A\).
Indeed, for \(q>0\), the definition of
\noderef{HierarchyStarredStage} identifies \(A=T_{n,q}\), and the first
line of (18) is false.  For \(q=0\), failure of (18) first says that
\(\beta\) is not missing on \(T_n\).  Equation (17) therefore gives
\(\beta\notin D_n\); equation (9) gives \(\alpha\notin D_n\).  Were
\(\beta\) missing on \(A\), the last alternative in (18) would hold,
again a contradiction.

In this complementary branch, exact prefixhood \(A\preccurlyeq B\),
the absence of a \(\beta\)-missing vertex on \(A\), and the choice of
\(v_\beta\) give
\[
  u\prec v_\beta\quad(u\in V(A)),\qquad
  v_\beta\prec y_{p-1}.                                    \tag{19}
\]
The first relation uses every vertex of \(A\), not merely its last
vertex.  Equation (16) gives
\(\beta\notin D_{n,q}\), the datum gives
\(\alpha\notin D_{n,q}\), and (16) also excludes \(\beta\) from each
individual displayed edge from the position of \(v_\beta\) up to the
position of \(y_{p-1}\).  Thus (19), the two stored-set exclusions, and
that exact unused-edge assertion are the additional hypotheses needed
later for the external stable-change provider.

The component in (20) is disjoint from the component rooted at \(y_p\).
If the first witnesses for \(\beta\) and \(\alpha\) coincide, both
colors are missing there, so expansion of
\noderef{TwoColorIncidenceGraph} makes that component an isolated
vertex, distinct from \(y_p\).  Otherwise, equality with the component
through \(y_p\) would put three distinct affected-color vertices of
incidence degree at most one on the finite spanning path, missing
respectively \(\beta,\alpha,\alpha\).  This contradicts
\noderef{TwoColorComponentSpanningPath} and
\noderef{TwoColorIncidenceDegreeAtMostTwo}.  The dichotomy in
\noderef{TwoColorComponentsDisjoint} gives disjointness, and the
spanning-path result supplies the possibly singleton
\noderef{TwoColorPathComponent} certificate at \(y_p\).

Define
\[
 \mu_3=\noderef{TwoColorComponentSwap}
       (G,\varphi_n,\beta,\alpha,\operatorname{inl}y_p).     \tag{21}
\]
We now keep the split on (18), rather than treating (17) as if it were
the first line of (18).

If \({\cal C}_{\beta\alpha}\) holds, apply
\noderef{StableChangeCasesOneTwo} with precisely (18), the elementary
prefix (3), the path certificate and component separation just proved,
and the six fields of the original minimum context.  Equations (10) and
(16) exclude both affected colors from every individual edge of
\(T-A\), so (21) fixes all those edges.  At a vertex of \(B^+\) outside
the swapped component no incident edge changes.  A vertex of \(B^+\)
inside it cannot be an endpoint other than one of the witnesses in the
disjoint component (20); at an internal vertex properness gives one edge
of each affected color, and the swap exchanges them.  Hence every
missing set on \(B^+\) is unchanged.  Since both affected colors are
outside \(D_{n,q}\), every stored first-missing position is unchanged.
These are the pointwise hypotheses of
\noderef{StableChangeCasesOneTwoFullTail}, which returns weak stability,
a mod-coloring, an actual corresponding maximum ETT on the same exact
\(T\), a rebuilt base and good hierarchy, and same-good data.

If \({\cal C}_{\beta\alpha}\) fails, apply
\noderef{StableChangeCaseThree} directly with ordered colors
\((\beta,\alpha)\), root \(\operatorname{inl}y_p\), the elementary
prefix (3), (19), the two \(D_{n,q}\)-exclusions, the exact unused-edge
consequence of (16), and the six fields of the original context.  Its
statement directly supplies weak stability, a mod-coloring, an actual
\noderef{CorrespondingExtendedTashkinovTree} on the full original
\(T\), maximum property, a rebuilt base, a good hierarchy, and
\noderef{SameGoodHierarchy}.  No full-tail or split-tail reconstruction
is invoked in this branch.  Nevertheless, equations (10) and (16) show
directly from the definition of (21) that every displayed edge above
\(A\) is fixed, and the endpoint/internal-vertex argument from the
preceding paragraph again proves vertexwise equality of all missing sets
on \(B^+\).

Thus in either branch we have the same full-tree certificate, same-good
data, fixed displayed edges above \(A\), and missing-set equality on
\(B^+\).  Construct \(P_3\) field by field.  Component separation keeps
\(\beta\) missing at \(v_\beta\), while the swap at the
\(\alpha\)-missing endpoint \(y_p\) makes \(\beta\) missing at \(y_p\);
these vertices are distinct by (2), so the replacement hierarchy is not
elementary.  Its six context fields are as follows:
\begin{enumerate}
\item properness is \noderef{TwoColorComponentSwapProper};
\item (21) preserves `none', hence preserves the exact unique uncolored
      edge in \noderef{IsOneEdgeUncolored};
\item rung and ladder-array agreement of the corresponding ETT transports
      the universally quantified lower-rung six-pack;
\item same-good identifies the divider count and transports the
      fewer-divider induction field;
\item the terminal tree is literally \(T\), and the equal stages and
      divider positions make both coordinates of
      \noderef{HierarchyCounterexampleMeasure} equal before the old
      lexicographic-minimality field is applied;
\item weak stability identifies the missing sets on the terminal ladder
      tree and transports the \(2n+11\) estimate.
\end{enumerate}
These data form an actual
\noderef{HierarchyReplacementPresentation} \(P_3\).

We next derive equation (19) of the paper inside \(P_3\).  Missing-set
equality on \(B^+\) retains the \(\theta\)-missingness at
\(y_{p-1}\).  Apply
\noderef{HierarchyCounterexamplePenultimateElementary} to the six fields
just constructed.  The resulting elementarity, together with
\(A\preccurlyeq B\) and \(y_{p-1}\notin V(B)\), says that \(\theta\)
is not missing on \(A\).  The branch assumption, transported by
same-good, says \(\theta\notin D_{n,q}\).  If \(\theta\in D_n\), the
definition \(D_{n,q}=D_n\setminus\overline\mu_3(A)\) would put it in
\(D_{n,q}\), a contradiction.  Therefore
\[
                        \theta\notin D_n.                    \tag{22}
\]
Apply \noderef{HierarchyTailMissingColorExclusion} under \(P_3\), with
the exact last vertex \(y_{p-1}\) of \(B^+\).  Its only exception is a
raw stored divider-zero seed, which (22) excludes, and hence
\[
             \theta\notin\mu_3\langle B^+-A\rangle.          \tag{23}
\]

Before the second use of \noderef{ChangeLemma}, we verify that its two
colors are distinct from the exact last edge.  The edge \(e_p\) is one
of the displayed edges of \(T-A\), so (16) gives
\(\varphi_n(e_p)\ne\beta\).  The datum gives
\(\varphi_n(e_p)=\theta\); consequently
\[
                         \beta\ne\theta.                    \tag{24}
\]
Similarly (10) gives \(\alpha\ne\theta\).  Thus the first swap affects
no \(\theta\)-colored edge, and in particular
\(\mu_3(e_p)=\theta\).  Moreover, \(\beta\) is present at
\(y_{p-1}\): it is missing at the distinct vertex
\(v_\beta\in V(B)\), (3) forbids a second missing witness in \(B^+\),
and the first swap preserves every missing set on \(B^+\).

Let \({\cal C}^{3}_{\beta\theta}\) be the exact Cases 1/2 proposition
for the coloring and presentation \(P_3\):
\[
 \begin{split}
 &(q>0\ \hbox{ and }
       \beta\in\overline\mu_3(T_{n,q}))\quad\hbox{or}\\
 &(q=0\ \hbox{ and }\ [\,
       \beta\in\overline\mu_3(T_n)\ \hbox{ or }
       (\beta\in\overline\mu_3(A),
        \beta,\theta\notin D_n^{3})\,]),                    \tag{25}
 \end{split}
\]
where \(D_n^{3}\) is the terminal stored set computed in \(P_3\).
The corresponding-ETT certificate identifies every raw stored array;
weak stability identifies the terminal-ladder missing sets; and
missing-set equality on \(B^+\) identifies missingness on
\(T_{n,q}\), \(T_n\), and \(A\).  Unfolding
\noderef{TashkinovStoredColors} therefore identifies \(D_n^3=D_n\).
Split on (25).  The following identifications show that this is the
same split retained from (18).  When \(q>0\), both predicates say only
that \(\beta\) is missing on \(T_{n,q}\), and missing-set equality on
\(B^+\) preserves exactly that assertion.  When \(q=0\), missingness on
\(T_n\) and \(A\) is likewise preserved.  If \(\beta\) is not missing
on \(T_n\), (17) gives \(\beta\notin D_n\); equation (9) gives
\(\alpha\notin D_n\), and (22) gives \(\theta\notin D_n^3=D_n\).
Consequently, in this remaining zero-divider alternative, the last
conjunct in (18) and the last conjunct in (25) are both equivalent to
the single assertion that \(\beta\) is missing on \(A\).  Hence (25)
holds exactly in the branch \({\cal C}_{\beta\alpha}\) and fails in its
complement.  The exact Cases 1/2 versus external classification is
therefore retained for the second recoloring.

Apply \noderef{ChangeLemma} to \(P_3\) with ordered colors
\((\beta,\theta)\) and the distinctness (24).  In the first branch,
(25) is exactly the Cases 1/2 alternative required by that lemma.
In the complementary branch, missing-set equality on \(B^+\) preserves
every assertion in (19): every vertex of \(A\)
precedes the first \(\beta\)-missing witness, and that witness precedes
the first \(\theta\)-missing vertex \(y_{p-1}\).  Same-good transports
\(\beta\notin D_{n,q}\), the branch assumption transports
\(\theta\notin D_{n,q}\), and the first swap fixes every displayed edge
above \(A\); hence \(\beta\) is still unused on the exact interval
between those witnesses.  For \(q>0\), the two transported
\(D_{n,q}\)-exclusions give the second alternative of the change lemma;
for \(q=0\), failure of (25), the retained failure of (18), and (22)
give \(\beta,\theta\notin D_n^3\), again its second alternative.  Thus
in either branch the change lemma gives
\[
 K^{\mu_3}_{\beta,\theta}(v_\beta)
      =K^{\mu_3}_{\beta,\theta}(y_{p-1}).                    \tag{26}
\]
If the first witnesses coincide, the common component is isolated.  If
they do not, equality with the component at \(y_p\) would create three
degree-at-most-one vertices on the spanning path, missing
\(\beta,\theta,\beta\).  The same three two-color component results used
after (20) show that the component through \(y_p\) is disjoint from
(26), and give its possibly singleton path certificate.  Set
\[
 \mu_4=\noderef{TwoColorComponentSwap}
       (G,\mu_3,\beta,\theta,\operatorname{inl}y_p).         \tag{27}
\]

Suppose first that (25) holds.  Apply
\noderef{StableChangeCasesOneTwo} relative to \(P_3\), using exactly
(25), the elementary penultimate prefix derived from its six context
fields, and the path and separation data above.  Use \(X=B^+\) in the
split-tail construction.  Equation (16), transported through the fixed
first swap, excludes \(\beta\) from every edge of \(T-A\), while (23)
excludes \(\theta\) from every edge of \(B^+-A\); (22) eliminates the
exceptional raw-stored seed.  Hence every individual prefix edge through
\(X\) is fixed by (27), and the endpoint/internal-vertex argument for
(26) gives missing-set equality throughout \(B^+\).

For the required route to \(X\), inspect the final PE, SE, or RE
\noderef{TashkinovSeriesStep} of the actual series in \(P_3\).
Restrict its closure route with
\noderef{TashkinovAugmentationSequenceRestrict}; weak stability protects
the PE source or the exact singleton SE/RE seed, and
\noderef{TashkinovAugmentationSequenceColoringTransport} transports the
route above \(A\) using the two pointwise equalities just proved.  Under
\(\mu_3\), the suffix \(B^+\to T\) is the one
\noderef{TashkinovAugmentation} adjoining the
\(\theta\)-colored edge \(e_p\), whose color is missing at
\(y_{p-1}\).

Both affected colors remain missing on \(X\), at witnesses in the
component disjoint from \(y_p\).  For stored first positions,
\(\beta\notin D_{n,q}\) by (16) and \(\theta\notin D_n\) by (22), so
every stored color is unaffected.  For the Gamma cutoff, (3) and
\(\theta\)'s missingness at \(y_{p-1}\) show that \(\theta\) is not
stage-missing and hence belongs to no last Gamma pair.  If \(\beta\) is
in no Gamma pair, the cutoff is vacuous.  Otherwise (16) supplies
\(\eta_h\in D_{n,q}\cap\overline\mu_3(B)\) with
\(\beta\in\Gamma^q_{\eta_h}\), so its first missing position is at most
\(|B|<|X|\).  These are all the quantified premises of
\noderef{StableChangeCasesOneTwoSplitTail}.  That theorem supplies weak
stability, a mod-coloring, the corresponding maximum ETT on the exact
full tree, a rebuilt good hierarchy, and same-good data under \(\mu_4\).

Suppose instead that (25) fails.  As established above, every vertex of
\(A\) precedes the first \(\beta\)-missing witness under \(\mu_3\),
that witness precedes \(y_{p-1}\), both \(\beta\) and \(\theta\) lie
outside the replacement \(D_{n,q}\), and \(\beta\) is unused on the
intervening exact edges.  Apply \noderef{StableChangeCaseThree} directly
to \(P_3\) with ordered colors \((\beta,\theta)\) and root
\(\operatorname{inl}y_p\).  It supplies weak stability, a mod-coloring,
an actual corresponding maximum ETT on the same full \(T\), a rebuilt
base and good hierarchy, and same-good data.  No split-tail
reconstruction is used in this external branch.  Component separation
and the endpoint/internal-vertex argument for (26) still show directly
that all missing sets on \(B^+\) are unchanged.

In both branches the exact edge-node \(e_p\) lies in the component rooted
at \(y_p\), because it is incident with \(y_p\) and has
\(\mu_3\)-color \(\theta\).  Expanding
\noderef{TwoColorIncidenceGraph} and (27) therefore gives
\(\mu_4(e_p)=\beta\).  Component separation keeps \(\theta\) missing at
\(y_{p-1}\), while properness makes \(e_p\) the unique incident
\(\theta\)-edge at \(y_p\); swapping it makes \(\theta\) missing at
\(y_p\).  The two vertices are distinct by (2), so the replacement
hierarchy is non-elementary.  Together with the provider's maximum ETT,
base, and good hierarchy, this constructs the one-step counterexample
over \(P_3\).

Build its minimum context with the same six separate arguments used for
\(P_3\): swap properness, preservation of `none', rung transport of the
lower six-pack, same-divider transport of fewer-divider induction,
equality of both measure coordinates before minimality, and
terminal-ladder missing-set equality before the large-set estimate.

The resulting correspondence is initially relative to \(P_3\), so make
\(P_4\) original-relative field by field.  Compose rung agreement and
each ladder-tree, coloring, stored-array, connecting-edge, chosen-vertex,
extension-kind, gamma, and delta equality from the two
\noderef{CorrespondingExtendedTashkinovTree} certificates.  At every
positive stable-coloring slot, compose the two stability relations by
the transitivity theorem \noderef{sc2}.  Compose the two hierarchy-stage
and divider-position equalities, the computed stored-set equalities, and
each Gamma-pair equality.  These are exactly the correspondence and
\noderef{SameGoodHierarchy} fields required by the original-relative
\noderef{HierarchyReplacementPresentation}; its context is the six-field
context just constructed.

Construct a genuine \noderef{HierarchySituationTwoCaseTwoData} for
\(P_4\).  Copy the exact tail indices, positions, non-incidence,
full-tail length, and path-number equation.  Its `.1' last-edge color is
\(\beta\), since \(\mu_4(e_p)=\beta\), and \(\beta\) is present at
\(y_{p-1}\) as proved before (25), with that fact preserved by the
missing-set equality on \(B^+\) in either provider branch.  Use
\(\theta\) as the Case 2 color:
it remains missing at \(y_{p-1}\), becomes missing at \(y_p\) under
(27), and is outside the final transported \(D_{n,q}\) by the original
nonstored-\(\theta\) branch assumption and the two same-good
equalities.  Thus every
field of the actual `.1' Case 2 datum is present.  Applying only the
already Sound \noderef{HierarchySituationTwoSubcaseTwoOne} to it and to
the actual minimum context of \(P_4\) gives the final contradiction.

Both alternatives for \(\theta\in D_{n,q}\) are impossible, completing
Subcase 2.2.
\end{proof}
